% RepNet Progress Update — Finalized Model Performance
% Same visual identity as final_presentation_modern.tex
% Compiles with pdfLaTeX.

\documentclass[10pt,aspectratio=169]{beamer}

% ====================
% Packages
% ====================
\usepackage[T1]{fontenc}
\usepackage{lmodern}
\usepackage{graphicx}
\usepackage{booktabs}
\usepackage{mathtools}
\usepackage{tikz}
\usepackage{hyperref}
\usepackage{pgfplots}
\pgfplotsset{compat=1.18}

% ====================
% UC Davis brand colors (digital palette)
% ====================
\definecolor{aggieblue}{HTML}{022851}
\definecolor{aggiegold}{HTML}{FFBF00}
\definecolor{aggieblueDark}{HTML}{011936}
\definecolor{aggiegoldDark}{HTML}{C99700}
\definecolor{coolgray}{RGB}{77,79,83}
\definecolor{boxgray}{RGB}{238,235,233}

% ====================
% Fonts
% ====================
\renewcommand{\familydefault}{\sfdefault}
\newcommand{\Raleway}{\sffamily}
\newcommand{\Lato}{\sffamily}

% ====================
% Beamer color slots
% ====================
\setbeamercolor*{palette primary}{fg=white,bg=aggieblue}
\setbeamercolor*{palette secondary}{fg=white,bg=coolgray}
\setbeamercolor*{palette tertiary}{fg=white,bg=aggieblueDark}
\setbeamercolor*{palette quaternary}{fg=aggieblue,bg=aggiegold}

\setbeamercolor{structure}{fg=aggieblue}
\setbeamercolor{titlelike}{fg=aggieblue}
\setbeamercolor{frametitle}{fg=white,bg=aggieblue}
\setbeamercolor{title}{fg=white,bg=aggieblue}
\setbeamercolor{author}{fg=aggieblue}
\setbeamercolor{institute}{fg=coolgray}
\setbeamercolor{date}{fg=coolgray}

\setbeamercolor{section in toc}{fg=aggieblue}
\setbeamercolor{subsection in toc}{fg=aggieblueDark}

% ====================
% Fonts (sizes/weights)
% ====================
\setbeamerfont{title}{family=\Raleway,series=\bfseries,size=\LARGE}
\setbeamerfont{subtitle}{family=\Lato,size=\normalsize}
\setbeamerfont{author}{family=\Lato,size=\small}
\setbeamerfont{institute}{family=\Lato,size=\footnotesize}
\setbeamerfont{date}{family=\Lato,size=\footnotesize}
\setbeamerfont{frametitle}{family=\Raleway,series=\bfseries,size=\large}
\setbeamerfont{section in toc}{family=\Raleway,series=\bfseries}
\setbeamerfont{subsection in toc}{family=\Lato,size=\small}
\setbeamerfont{footline}{family=\Lato,size=\scriptsize}

% ====================
% Frame title with Aggie Gold accent bar
% ====================
\setbeamertemplate{frametitle}{%
  \nointerlineskip
  \begin{beamercolorbox}[wd=\paperwidth,ht=3.0ex,dp=1.4ex,leftskip=1em,rightskip=1em]{frametitle}%
    \usebeamerfont{frametitle}\insertframetitle%
  \end{beamercolorbox}%
  \begin{beamercolorbox}[wd=\paperwidth,ht=0.35ex,dp=0pt]{palette quaternary}%
  \end{beamercolorbox}%
  \vspace{1.0ex}
}

\beamertemplatenavigationsymbolsempty

% ====================
% Table of Contents Customization
% ====================
\setbeamertemplate{section in toc}{\leavevmode\inserttocsection\par}

\setbeamertemplate{subsection in toc}{%
  \leavevmode\hspace{1.5em}%
  \color{aggiegoldDark}\raise0.4ex\hbox{\vrule width 0.4ex height 0.4ex}%
  \hspace{0.5em}\color{aggieblueDark}\inserttocsubsection\par
}

% ====================
% Footer
% ====================
\newcommand{\footercontent}[1]{\renewcommand{\insertfootercontent}{#1}}
\newcommand{\insertfootercontent}{}

\setbeamertemplate{footline}{%
  \ifx\insertfootercontent\@empty\else
    \begin{beamercolorbox}[wd=\paperwidth,ht=0.3ex,dp=0pt]{palette quaternary}\end{beamercolorbox}%
    \begin{beamercolorbox}[wd=\paperwidth,ht=2.4ex,dp=1.2ex,leftskip=1em,rightskip=1em]{palette primary}%
      \usebeamerfont{footline}\insertfootercontent\hfill\insertframenumber{}/\inserttotalframenumber%
    \end{beamercolorbox}%
  \fi
}

% ====================
% Title / authors
% ====================
\title{Risk Estimation of Preeclampsia}
\subtitle{RepNet}
\author{Ahmed Khan \and Lawrencedel Legaspi \and Selena Phu}
\institute{Department of Electrical and Computer Engineering \\ University of California, Davis}
\date{\today}

\titlegraphic{
\includegraphics[height=2.0cm]{logo.png}}

\footercontent{%
  \href{https://github.com/AhmedKhan-GH/repnet}{github.com/AhmedKhan-GH/repnet}
  \enspace$\bullet$\enspace
  UC Davis Senior Design 2026%
}

% ====================
% Title page template
% ====================
\setbeamertemplate{title page}{%
  \begin{tikzpicture}[remember picture,overlay]
    \fill[aggieblue]
      ([yshift=1cm]current page.north west)
      rectangle
      ([yshift=-3.8cm]current page.north east);
    \fill[aggiegold]
      ([yshift=-3.8cm]current page.north west)
      rectangle
      ([yshift=-3.95cm]current page.north east);
    \node[anchor=center,text=white,align=center,text width=14cm,
          font=\Raleway\bfseries]
      at ([yshift=-1.6cm]current page.north)
      {\usebeamerfont{title}\inserttitle};
    \node[anchor=center,text=aggiegold,align=center,text width=14cm,
          font=\Lato]
      at ([yshift=-2.4cm]current page.north)
      {\usebeamerfont{subtitle}\insertsubtitle};
    \node[anchor=north]
      at ([yshift=-4.3cm]current page.north)
      {\inserttitlegraphic};
    \node[anchor=north,text=aggieblue,align=center,font=\Lato]
      at ([yshift=-6.5cm]current page.north)
      {\usebeamerfont{author}\insertauthor};
    \node[anchor=north,text=coolgray,align=center,font=\Lato]
      at ([yshift=-7.1cm]current page.north)
      {\usebeamerfont{institute}\insertinstitute};
    \node[anchor=north,text=coolgray,align=center,font=\Lato]
      at ([yshift=-8.0cm]current page.north)
      {\usebeamerfont{date}\insertdate};
  \end{tikzpicture}%
}

% ====================
% Hidden section declarations to populate the TOC
% ====================
\section{Introduction}
\section{Decision-Tree Modeling}
\subsection{Dataset Profile}
\subsection{Feature Extraction}
\subsection{Feature Importance}
\subsection{Model Results}
\section{Neural Network Architecture}
\section{Conclusion}
\section{Future Work}


% ====================
% Document
% ====================
\begin{document}

% ---------- Title Slide ----------
{%
  \setbeamertemplate{footline}{}%
  \begin{frame}[plain,noframenumbering]
    \titlepage
  \end{frame}
}

% ---------- Table of Contents Slide ----------
\begin{frame}{Contents}
  \tableofcontents
\end{frame}

% ---------- Transition Slide ----------
{%
  \setbeamertemplate{footline}{}%
  \begin{frame}[plain,noframenumbering]
    \begin{tikzpicture}[remember picture,overlay]
      \fill[aggieblue] (current page.south west) rectangle (current page.north east);
    \end{tikzpicture}
    \vfill
    \centering
    {\Raleway\bfseries\Huge\color{white} Introduction}\\[0.6em]
    {\color{aggiegold}\rule{6cm}{0.06cm}}
    \vfill
  \end{frame}
}

\begin{frame}{\textbf{Introduction: Preeclampsia Detection via RepNet}}
    
    % Section 1: The Clinical Problem
    \textbf{Clinical Motivation}
    \begin{itemize}
        \item \textbf{Preeclampsia:} A severe hypertensive pregnancy disorder that can lead to maternal and fetal mortality.
        \item \textbf{Complications:} Liver damage, low platelet count, and preterm birth.
        \item \textbf{Risk of Eclampsia:} Can rapidly progress to affect brain function, putting the mother at risk of seizures or coma.
    \end{itemize}
    
    \vspace{0.3cm} % Clean spacing between sections
    
    % Section 2: The Proposed Machine Learning Solution
    \textbf{Our Approach: RepNet}
    \begin{itemize}
        \item A pair of machine learning models combining 12-lead ECGs and clinical data to estimate preeclampsia risk:
        \begin{itemize}
            \item \textbf{RepNet-RF:} Decision-tree modeling trained on engineered features and clinical data.
            \item \textbf{RepNet-SE:} A depth-wise separable CNN with squeeze-excitation layers and cross-lead attention mechanisms trained on raw ECG waveforms.
        \end{itemize}
        \item \textbf{Dataset:} Evaluate on 2,191 clinical 12-lead ECGs (337 PreE, 15.4\%) across 1,383 patients at UC Davis Health.
    \end{itemize}

\end{frame}

% ---------- Transition Slide ----------
{%
  \setbeamertemplate{footline}{}%
  \begin{frame}[plain,noframenumbering]
    \begin{tikzpicture}[remember picture,overlay]
      \fill[aggieblue] (current page.south west) rectangle (current page.north east);
    \end{tikzpicture}
    \vfill
    \centering
    {\Raleway\bfseries\Huge\color{white} Decision-Tree Modeling}\\[0.6em]
    {\color{aggiegold}\rule{6cm}{0.06cm}}
    \vfill
  \end{frame}
}

% ---------- Subsection Slide: Dataset Profile ----------
\begin{frame}{Decision-Tree Modeling --- Dataset Profile}
  \begin{columns}[T]
    
    % Left Column: Cohort Breakdown Table & Details
    \begin{column}{0.50\textwidth}
      \centering
      {\small\Raleway\bfseries\color{aggieblue} ECG Dataset}\\[0.4em]
      \resizebox{\textwidth}{!}{%
        \begin{tabular}{lccc}
          \toprule
          \textbf{Dataset} & \textbf{\shortstack{Original Total\\(Preeclamptic/Normal)}} & \textbf{\shortstack{Dropped ECGs\\(Preeclamptic/Normal)}} & \textbf{\shortstack{Final Total\\(Preeclamptic/Normal)}} \\
          \midrule
          Unbalanced & 337 / 1854 & 8 / 24 & 329 / 1830 \\
          Balanced   & 187 / 182  & 6 / 1  & 181 / 181 \\
          \bottomrule
        \end{tabular}
      }%
      
      \vspace{0.8em}
      \begin{flushleft}
        \usebeamerfont{footline}\color{coolgray}
        \textbf{MUSE Diagnostics:} This dataset includes a structured metadata CSV containing \textbf{656 distinct features}.
        
        \vspace{0.4em}
        \textbf{Exclusion Criteria:} Samples dropped due to:
        \begin{itemize}
          \item Delineation issues using certain cleaning methods for each lead within \textbf{NeuroKit} framework.
          \item Empty data within a lead.
        \end{itemize}
      \end{flushleft}
    \end{column}
    
    % Right Column: Preprocessed Lead Signal Image
    \begin{column}{0.46\textwidth}
      \centering
      {\small\Raleway\bfseries\color{aggieblue} Lead Visualizations}\\[0.4em]
      \begin{beamercolorbox}[wd=\textwidth,rounded=true,shadow=false,leftskip=0.2em,rightskip=0.2em,topskip=0.2em,bottomskip=0.2em]{boxgray}
        \centering
        \includegraphics[width=0.96\textwidth,height=3.8cm,keepaspectratio]{4leadsexample.png}
      \end{beamercolorbox}
      \vspace{0.3em}
      {\scriptsize\Lato\color{coolgray} Representative trace behavior of select isolated signal leads (e.g., I, II, III, aVR) after preprocessing filters are applied.}
    \end{column}

  \end{columns}
\end{frame}

% ---------- Subsection Slide: Feature Extraction ----------
\begin{frame}{Decision-Tree Modeling --- Feature Extraction}
  \begin{columns}[T]
    
    % Left Column: Neurokit Feature Pipeline
    \begin{column}{0.48\textwidth}
      \centering
      {\small\Raleway\bfseries\color{aggieblue} Neurokit ECG Features}\\[0.4em]
      \resizebox{\textwidth}{!}{%
        \begin{tabular}{ll}
          \toprule
          \textbf{Feature} & \textbf{Equation} \\
          \midrule
          $P_{\max}$ & $\max_{\ell \in \mathcal{L}}(P_{\text{off}}^{\ell} - P_{\text{on}}^{\ell})$ \\
          $P_{\min}$ & $\min_{\ell \in \mathcal{L}}(P_{\text{off}}^{\ell} - P_{\text{on}}^{\ell})$ \\
          P-Wave Disp. & $P_{\max} - P_{\min}$ \\
          QT Dispersion & $QT_{\max} - QT_{\min}$ \\
          $T_{pe}$ Interval & $\text{mean}(T_{\text{off}} - T_{\text{peak}})$ \\
          $T_{pe}/\text{QT}$ & $T_{pe} \text{ Interval} / \text{mean}(T_{\text{off}} - R_{\text{on}})$ \\
          $T_{pe}/\text{QTc}$ & $(T_{pe}/\text{QT}) / \sqrt{\text{mean}(\Delta R_{\text{peak}})}$ \\
          ST Elevation & $\boldsymbol{1} \text{ if } > 50\%\text{ elevated, } \boldsymbol{0}\text{ otherwise}$ \\
          ST Depression & $\boldsymbol{1} \text{ if } > 50\%\text{ depressed, } \boldsymbol{0}\text{ otherwise}$ \\
          Avg. Amp. Diff. & $\text{mean}(A_{\text{ST}} - A_{\text{PR}})$ \\
          Med. Amp. Diff. & $\text{med}(A_{\text{ST}} - A_{\text{PR}})$ \\
          Max. Amp. Diff. & $\max(A_{\text{ST}} - A_{\text{PR}})$ \\
          Min. Amp. Diff. & $\min(A_{\text{ST}} - A_{\text{PR}})$ \\
          Frontal QRS-T & $|\text{QRS Axis} - \text{T Axis}|$ \\
          \bottomrule
        \end{tabular}
      }%
    \end{column}
    
    % Right Column: MUSE Feature Pipeline + Footnote Text (Clean & Compact)
    \begin{column}{0.48\textwidth}
      \centering
      {\small\Raleway\bfseries\color{aggieblue} MUSE-derived Features}\\[0.4em]
      \resizebox{\textwidth}{!}{%
        \begin{tabular}{ll}
          \toprule
          \textbf{Feature} & \textbf{Equation} \\
          \midrule
          $P_{\max}$ & $\max(\text{P\_Duration for 12 leads})$ \\
          $P_{\min}$ & $\min(\text{P\_Duration for 12 leads})$ \\
          P-Wave Disp. & $P_{\max} - P_{\min}$ \\
          QT Dispersion & $QT = \sum \text{comp. durations} + \frac{1}{8}RR; \, \max(QT) - \min(QT)$ \\
          $T_{pe}$ Interval & $\text{T\_Duration} + \text{T\_PrimeDuration} - \text{STEtoTPeak}$ \\
          $T_{pe}/\text{QT}$ & $T_{pe} \text{ Interval} / \text{QT}$ \\
          $T_{pe}/\text{QTc}$ & $T_{pe}/\text{QT} \cdot \sqrt{RR}$ \\
          ST Elevation & $(\text{BITFLGS} \ \& \ 8) \neq 0$ \\
          ST Depression & $(\text{BITFLGS} \ \& \ 4) \neq 0$ \\
          QRS Axis & $\text{RAxis from metadata}$ \\
          T Axis & $\text{TAxis from metadata}$ \\
          Frontal QRS-T & $|\text{QRS Axis} - \text{T Axis}|$ \\
          \bottomrule
        \end{tabular}
      }%
      
      \vspace{1.2em}
      \begin{beamercolorbox}[wd=\textwidth,rounded=true,shadow=false,leftskip=0.6em,rightskip=0.6em,topskip=0.4em,bottomskip=0.4em]{boxgray}
        \usebeamerfont{footline}\color{coolgray}
        \textbf{Implementation Note:} We extracted QRS and T Axis in Neurokit, though we do not have a definite equation.
      \end{beamercolorbox}
    \end{column}

  \end{columns}
\end{frame}

% ---------- Subsection Slide: Model Performance ----------
\begin{frame}{Decision-Tree Modeling --- Feature Importance}
  \centering
  {\small\Raleway\bfseries\color{aggieblue} Random Forest Performance}\\[0.3em]
  
  \resizebox{0.95\textwidth}{!}{%
    \begin{tabular}{lcccccccc}
      \toprule
      \textbf{Method} & \textbf{\# Features} & \textbf{AUROC} & \textbf{AUPRC} & \textbf{Specificity} & \textbf{Sensitivity} & \textbf{Bal. Accuracy} & \textbf{Precision} & \textbf{F1-Score} \\
      \midrule
      Original & 723 & $69.93 \pm 3.30$ & $32.55 \pm 3.71$ & $73.74 \pm 3.26$ & $56.12 \pm 8.60$ & $64.93 \pm 3.71$ & $27.97 \pm 2.62$ & $37.20 \pm 3.89$ \\
      Pearson Corr. & 363 & $70.97 \pm 3.86$ & $35.14 \pm 7.29$ & $74.57 \pm 4.37$ & $54.89 \pm 3.83$ & $64.73 \pm 2.62$ & $28.64 \pm 4.02$ & $37.45 \pm 3.48$ \\
      \textbf{ANOVA} & \textbf{191} & $\mathbf{74.54 \pm 3.88}$ & $\mathbf{37.71 \pm 5.03}$ & $74.15 \pm 3.08$ & $\mathbf{62.04 \pm 10.59}$ & $\mathbf{68.10 \pm 4.37}$ & $\mathbf{30.08 \pm 2.62}$ & $\mathbf{40.36 \pm 4.10}$ \\
      RFECV & 175 & $73.31 \pm 3.77$ & $38.12 \pm 6.08$ & $74.48 \pm 2.73$ & $59.32 \pm 10.49$ & $66.90 \pm 3.96$ & $29.30 \pm 1.61$ & $39.08 \pm 3.60$ \\
      \midrule
      Intersection & 176 & $73.09 \pm 4.43$ & $36.58 \pm 7.34$ & $73.06 \pm 4.90$ & $61.15 \pm 14.82$ & $71.24 \pm 2.02$ & $28.75 \pm 1.54$ & $38.76 \pm 4.25$ \\
      \bottomrule
    \end{tabular}
  }%

  \vspace{0.6em}
  \begin{beamercolorbox}[wd=0.95\textwidth,rounded=true,leftskip=0.8em,rightskip=0.8em,topskip=0.4em,bottomskip=0.4em]{boxgray}
    \usebeamerfont{footline}\color{coolgray}
    \textbf{Performance Insight:} Features chosen \textbf{sequentially} yielded stronger performance results compared to doing the importances \textbf{individually} and intersecting the common features.
  \end{beamercolorbox}
\end{frame}

% ---------- Subsection Slide: Feature Importance Breakdown ----------
\begin{frame}{Decision-Tree Modeling --- Feature Importance}
  \begin{columns}[T]
    
    % Left Column: Methodological Framework details
    \begin{column}{0.42\textwidth}
      \usebeamerfont{footline}\color{coolgray}
      \textbf{Cross-Validation Strategy:}
      \begin{itemize}
        \item\textbf{StratifiedGroupKFold} with 5 folds grouped by the Patient's MRN
        \item\textbf{RandomForest} with \textbf{RandomUnderSampler} inside the folds
      \end{itemize}

      \vspace{0.2em}
      \textbf{Feature Selection:}
      \begin{itemize}
        \item \textbf{Pearson Correlation:} Drops one feature out of colinear pairs sharing a correlation value $> 0.85$.
        \item \textbf{ANOVA (Analysis of Variance):} Keeps features with $p < 0.05$.
        \item \textbf{RFECV:} Used roc\_auc scoring to recursively remove features
      \end{itemize}
      
      \vspace{0.2em}
      \begin{beamercolorbox}[wd=\textwidth,rounded=true,leftskip=0.5em,rightskip=0.5em,topskip=0.3em,bottomskip=0.3em]{boxgray}
        \scriptsize\Lato\color{coolgray}
        \textbf{Hindsight Note:} Evaluating RFECV on \textit{average\_precision} or \textit{balanced\_accuracy} may have been better for the primary goal of detecting preeclampsia.
      \end{beamercolorbox}
    \end{column}
    
    % Right Column: Lead distribution chart
    \begin{column}{0.54\textwidth}
      \centering
      {\small\Raleway\bfseries\color{aggieblue} Feature Counts Across Leads}\\[0.3em]
      \begin{beamercolorbox}[wd=\textwidth,rounded=true,leftskip=0.1em,rightskip=0.1em,topskip=0.1em,bottomskip=0.1em]{boxgray}
        \centering
        \includegraphics[width=0.98\textwidth,height=4.2cm,keepaspectratio]{leadimportance.png}
      \end{beamercolorbox}
      \vspace{0.4em}
      \begin{flushleft}
        \usebeamerfont{footline}\color{coolgray}
        \textbf{Lead Significance:} Graph maps feature distribution counts across different leads, showing which leads may be more important.
      \end{flushleft}
    \end{column}

  \end{columns}
\end{frame}

% ---------- Subsection Slide: Feature Selection Matrix ----------
\begin{frame}{Decision-Tree Modeling --- Feature Heatmap}
  
  % [onlytextwidth] removes outer margins to give the wide heatmap maximum breathing room
  \begin{columns}[c, onlytextwidth] 
    
    % Left Column: Maximized Heatmap Allocation (82% width)
    \begin{column}{0.82\textwidth}
      \centering
      {\small\Raleway\bfseries\color{aggieblue} Feature Selection by Methods --- 12 Leads}\\[0.3em]
      \begin{beamercolorbox}[wd=\textwidth,rounded=true,leftskip=0.1em,rightskip=0.1em,topskip=0.1em,bottomskip=0.1em]{boxgray}
        \centering
        \includegraphics[width=\textwidth,height=4.2cm,keepaspectratio]{Heatmap12lead2.png}
      \end{beamercolorbox}
    \end{column}
    
    % Right Column: Downsized Legend Allocation (15% width)
    \begin{column}{0.15\textwidth}
      \centering
      \vspace{1.2em} % Subtle vertical alignment adjustment
      \begin{beamercolorbox}[wd=\textwidth,rounded=true,leftskip=0.1em,rightskip=0.1em,topskip=0.1em,bottomskip=0.1em]{boxgray}
        \centering
        \includegraphics[width=\textwidth,keepaspectratio]{heatmaplegend.png}
      \end{beamercolorbox}
    \end{column}

  \end{columns}

  \vspace{0.6em}
  \hr
  
  % Contextual Breakdown spanning beneath the images
  \begin{columns}[T, onlytextwidth]
    \begin{column}{0.48\textwidth}
      \usebeamerfont{footline}\color{coolgray}
      \textbf{Feature Selection Methods:} Previous Pearson, ANOVA, and RFECV (roc\_auc), but now includes RFECV (balanced\_accuracy and average\_precision)
    \end{column}
    \begin{column}{0.48\textwidth}
      \usebeamerfont{footline}\color{coolgray}
      \textbf{Pipeline Constraint Note:} The additional balanced\_accuracy and average\_precision are not included for the future ensemble modeling as these were included afterwards. This heatmap is just a showing of how important overall features may be and can tell us whether a feature is universally important across leads.
    \end{column}
  \end{columns}
\end{frame}

% ---------- Subsection Slide: Final Sequential Feature Importance ----------
\begin{frame}{Decision-Tree Modeling --- Final Feature Importance}
  \begin{columns}[T]
    
    % Left Column: Deep Dive into the exact features you called out
    \begin{column}{0.44\textwidth}
      \usebeamerfont{footline}\color{coolgray}
      \textbf{Key Predictor Breakdown:}
      \begin{itemize}
        \item \textbf{Highest seen Features:} \texttt{BITFLGS}, \texttt{QRS Intrinsicoid}, and \texttt{STEtoTPeak} were most frequent across the 12-lead array.
        \item \textbf{NeuroKit Local Features:} The most critical lead-isolated extractors were \textbf{Tp-e}, \textbf{Tp-e/QT}, and \textbf{ST Depression}.
        \item \textbf{Surviving Global Features:} Beyond individual leads, the non-isolated features that cleared selection were \textbf{$P_{\max}$}, \textbf{$P_{\min}$}, \textbf{T Axis}, and the \textbf{Frontal QRS-T Angle}.
      \end{itemize}
    \end{column}
    
    % Right Column: The Sequential Feature Importance Bar Chart
    \begin{column}{0.54\textwidth}
      \centering
      {\small\Raleway\bfseries\color{aggieblue} Final Active Modeling Features}\\[0.3em]
      \begin{beamercolorbox}[wd=\textwidth,rounded=true,leftskip=0.1em,rightskip=0.1em,topskip=0.1em,bottomskip=0.1em]{boxgray}
        \centering
        % Generates the horizontal bar chart mapping counts/importance from sequential selection
        \includegraphics[width=0.98\textwidth,height=3.8cm,keepaspectratio]{featureimportance.png}
      \end{beamercolorbox}
      \vspace{0.3em}
      \begin{flushleft}
        \usebeamerfont{footline}\color{coolgray}
        \textbf This distribution shows the final feature space pruned through the Pearson $\rightarrow$ ANOVA $\rightarrow$ RFECV sequential chain.
      \end{flushleft}
    \end{column}

  \end{columns}
\end{frame}

% ---------- Subsection Slide: Model Performance Comparison ----------
\begin{frame}{Decision-Tree Modeling --- Individual vs. Ensemble Performance}
  
  % Tighten padding slightly to give the text boxes maximum horizontal clearance
  \setlength{\tabcolsep}{3.5pt}
  
  \begin{columns}[T, onlytextwidth]
    
    % Left Column: Custom Commentary Text Box Placeholder (34% width)
    \begin{column}{0.34\textwidth}
      \vspace{0.2em}
      \begin{beamercolorbox}[wd=\textwidth,rounded=true,leftskip=0.8em,rightskip=0.8em,topskip=0.8em,bottomskip=0.8em]{boxgray}
        
        \usebeamerfont{footline}\color{coolgray}
        
     
        \textbf{Model:}
\begin{itemize}
    \item \textbf{Unweighted Voting+} consists of: RandomForest, LightGBM, XGBoost, CatBoost, LogisticRegression, and SVM.
    \item \textbf{Unweighted/Weighted Voting} consists of: RandomForest, LightGBM, XGBoost, CatBoost, and AdaBoost.
    \item \textbf{Weights:} Votes are weighted [4, 2, 3, 4, 1] (respective to the models) based on AUPRC and sensitivity.
    \item Each individual model is tuned on the 175 features using RandomizedSearchCV.
\end{itemize}
        
      \end{beamercolorbox}
    \end{column}
    
    % Right Column: Fixed Native LaTeX Performance Tables (63% width)
    \begin{column}{0.63\textwidth}
      \centering
      
      % Top Section: Standalone Base Models Table
      {\scriptsize\Raleway\bfseries\color{aggieblue} Standalone Tree-Based Classifiers}\\[0.3em]
      \begin{beamercolorbox}[wd=\textwidth,rounded=true,leftskip=0.2em,rightskip=0.2em,topskip=0.4em,bottomskip=0.4em]{boxgray}
        \tiny
        % \resizebox ensures the table never clips past the gray container boundaries
        \resizebox{\textwidth}{!}{
          \begin{tabular}{lcccccc}
            \textbf{Model} & \textbf{AUROC} & \textbf{AUPRC} & \textbf{Specificity} & \textbf{Sensitivity} & \textbf{Balanced Accuracy} & \textbf{F1-Score} \\
            \hline
            RandomForest & $74.49 \pm 3.77$ & $42.56 \pm 7.60$ & $74.37 \pm 3.20$ & $61.43 \pm 11.36$ & $67.90 \pm 4.30$ & $40.10 \pm 3.86$ \\
            LightGBM     & $73.25 \pm 3.47$ & $37.30 \pm 7.47$ & $72.95 \pm 4.33$ & $59.62 \pm 9.49$  & $66.28 \pm 2.62$ & $38.24 \pm 1.88$ \\
            XGBoost      & $74.84 \pm 4.50$ & $40.76 \pm 7.14$ & $74.26 \pm 3.76$ & $61.44 \pm 11.11$ & $67.85 \pm 4.28$ & $40.09 \pm 3.93$ \\
            CatBoost     & $74.55 \pm 4.94$ & $42.24 \pm 8.67$ & $75.96 \pm 4.81$ & $60.53 \pm 13.37$ & $68.24 \pm 4.56$ & $41.70 \pm 3.90$ \\
            AdaBoost     & $69.86 \pm 4.43$ & $33.91 \pm 7.71$ & $67.92 \pm 7.96$ & $59.90 \pm 13.17$ & $63.91 \pm 3.04$ & $35.14 \pm 2.44$ \\
          \end{tabular}
        }
      \end{beamercolorbox}
      
      \vspace{0.6em}
      
      % Bottom Section: Ensemble Architectures Table
      {\scriptsize\Raleway\bfseries\color{aggieblue} Multi-Model Ensemble Frameworks}\\[0.3em]
      \begin{beamercolorbox}[wd=\textwidth,rounded=true,leftskip=0.2em,rightskip=0.2em,topskip=0.4em,bottomskip=0.4em]{boxgray}
        \tiny
        \resizebox{\textwidth}{!}{
          \begin{tabular}{lcccccc}
            \textbf{Ensemble} & \textbf{AUROC} & \textbf{AUPRC} & \textbf{Specificity} & \textbf{Sensitivity} & \textbf{Balanced Accuracy} & \textbf{F1-Score} \\
            \hline
            Unweighted Voting  & $74.54 \pm 4.16$ & $41.81 \pm 7.86$ & $74.26 \pm 3.95$ & $59.32 \pm 10.70$ & $66.79 \pm 3.60$ & $38.94 \pm 2.99$ \\
            Weighted Voting    & $74.75 \pm 4.35$ & $42.26 \pm 7.92$ & $74.75 \pm 4.48$ & $60.23 \pm 10.88$ & $67.49 \pm 3.45$ & $39.80 \pm 2.76$ \\
            \textbf{Unweighted Voting+} & $\mathbf{75.72 \pm 3.94}$ & $41.58 \pm 8.03$ & $74.64 \pm 4.70$ & $\mathbf{62.35 \pm 10.93}$ & $\mathbf{68.50 \pm 3.22}$ & $40.85 \pm 2.24$ \\
            Stacking           & $74.77 \pm 4.35$ & $42.43 \pm 7.59$ & $74.37 \pm 4.18$ & $59.62 \pm 10.80$ & $67.00 \pm 3.46$ & $39.17 \pm 2.79$ \\
            Bagging CatBoost   & $74.81 \pm 4.98$ & $41.96 \pm 8.86$ & $74.70 \pm 4.68$ & $59.62 \pm 15.19$ & $67.16 \pm 5.59$ & $39.10 \pm 5.21$ \\
          \end{tabular}
        }
      \end{beamercolorbox}
      
    \end{column}

  \end{columns}
\end{frame}

% ---------- Transition Slide ----------
{%
  \setbeamertemplate{footline}{}%
  \begin{frame}[plain,noframenumbering]
    \begin{tikzpicture}[remember picture,overlay]
      \fill[aggieblue] (current page.south west) rectangle (current page.north east);
    \end{tikzpicture}
    \vfill
    \centering
    {\Raleway\bfseries\Huge\color{white} Neural Network Architecture}\\[0.6em]
    {\color{aggiegold}\rule{6cm}{0.06cm}}
    \vfill
  \end{frame}
}

% ---------- Transition Slide ----------
{%
  \setbeamertemplate{footline}{}%
  \begin{frame}[plain,noframenumbering]
    \begin{tikzpicture}[remember picture,overlay]
      \fill[aggieblue] (current page.south west) rectangle (current page.north east);
    \end{tikzpicture}
    \vfill
    \centering
    {\Raleway\bfseries\Huge\color{white} Conclusion}\\[0.6em]
    {\color{aggiegold}\rule{6cm}{0.06cm}}
    \vfill
  \end{frame}
}

% ---------- Transition Slide ----------
{%
  \setbeamertemplate{footline}{}%
  \begin{frame}[plain,noframenumbering]
    \begin{tikzpicture}[remember picture,overlay]
      \fill[aggieblue] (current page.south west) rectangle (current page.north east);
    \end{tikzpicture}
    \vfill
    \centering
    {\Raleway\bfseries\Huge\color{white} Future Work}\\[0.6em]
    {\color{aggiegold}\rule{6cm}{0.06cm}}
    \vfill
  \end{frame}
}

% ---------- Subsection Slide: Future Directions & Limitations ----------
\begin{frame}{Future Work}
  \begin{columns}[T, onlytextwidth]
    
    % Left Column: Operational Extensions
    \begin{column}{0.48\textwidth}
      {\small\Raleway\bfseries\color{aggieblue} Scaling \& Translation}\\[0.4em]
      
      \begin{beamercolorbox}[wd=\textwidth,rounded=true,leftskip=0.6em,rightskip=0.6em,topskip=0.5em,bottomskip=0.5em]{boxgray}
        \usebeamerfont{footline}\color{coolgray}
        \textbf{Lead Reduction \& Wearables:}
        \begin{itemize}
          \item Reduced lead sets (e.g., leads I, II, V2, V4) or single-lead models could enable screening via smart wearables such as Apple Watch ECGs.
        \end{itemize}
      \end{beamercolorbox}
      
      \vspace{0.5em}
      
      \begin{beamercolorbox}[wd=\textwidth,rounded=true,leftskip=0.6em,rightskip=0.6em,topskip=0.5em,bottomskip=0.5em]{boxgray}
        \usebeamerfont{footline}\color{coolgray}
        \textbf{Transfer Learning:}
        \begin{itemize}
          \item Pretraining on large public ECG corpora (PTB-XL, MIMIC-IV-ECG) or leveraging ECG foundation models before fine-tuning on the UCDH cohort.
        \end{itemize}
      \end{beamercolorbox}
    \end{column}
    
    % Right Column: Robustness & Contextual Inputs
    \begin{column}{0.48\textwidth}
      {\small\Raleway\bfseries\color{aggieblue} Generalizability \& Context}\\[0.4em]
      
      \begin{beamercolorbox}[wd=\textwidth,rounded=true,leftskip=0.6em,rightskip=0.6em,topskip=0.5em,bottomskip=0.5em]{boxgray}
        \usebeamerfont{footline}\color{coolgray}
        \textbf{Multi-Site Validation:}
        \begin{itemize}
          \item Current results are from a single site; external validation on diverse populations is needed to assess generalizability.
        \end{itemize}
      \end{beamercolorbox}
      
      \vspace{0.5em}
      
      \begin{beamercolorbox}[wd=\textwidth,rounded=true,leftskip=0.6em,rightskip=0.6em,topskip=0.5em,bottomskip=0.5em]{boxgray}
        \usebeamerfont{footline}\color{coolgray}
        \textbf{Demographic Integration:}
        \begin{itemize}
          \item Incorporating maternal age, gestational age, and clinical history as auxiliary inputs alongside ECG data.
        \end{itemize}
      \end{beamercolorbox}
    \end{column}

  \end{columns}
\end{frame}



% ---------- Closing Slide ----------
{%
  \setbeamertemplate{footline}{}%
  \begin{frame}[plain,noframenumbering]
    \begin{tikzpicture}[remember picture,overlay]
      \fill[aggieblue] (current page.south west) rectangle (current page.north east);
    \end{tikzpicture}
    \vfill
    \centering
    {\Raleway\bfseries\Huge\color{white} Thank you}\\[0.6em]
    {\color{aggiegold}\rule{4cm}{0.06cm}}\\[0.8em]
    {\Lato\large\color{white} Questions?}
    \vfill
  \end{frame}
}

\end{document}